\documentclass[10pt,a4paper]{amsart}
\usepackage[margin=19mm]{geometry}
\usepackage{amsmath,amssymb,mathtools}
\usepackage[scr=rsfs,cal=euler]{mathalfa}
\usepackage{xfrac}
\usepackage[unicode,hidelinks,bookmarks=false]{hyperref}
\newcommand{\ie}{i.e.\ }
\newcommand{\cf}{cf.\ }
\newcommand{\resp}{respectively\ }
\newcommand{\Subsection}{\subsection}
\providecommand{\nobreakdash}{\nobreak\hbox{-}}
\setlength{\emergencystretch}{2em}
\multlinegap0pt
\usepackage{nicefrac}
\usepackage{booktabs}
\usepackage{multirow}
\usepackage{xcolor}
\usepackage{float}
\usepackage{mathtools}
\usepackage{tcolorbox}
\usepackage{braket}
\usepackage{tikz}
\usepackage{amssymb}
\usepackage[shortlabels]{enumitem}
\usepackage{algorithm}\usepackage{algpseudocode}
\newtheorem{theorem}{Theorem}
\newcommand{\GF}{\mathrm{GF}}
\newcommand{\cL}{\mathcal{L}}
\newcommand{\defeq}{:=}
\title{LLM-Text Watermarking based on Lagrange Interpolation}
\author{Jaros{\l}aw Janas, Pawe{\l} Morawiecki, Josef Pieprzyk}
\date{Source: arXiv:2505.05712v4 (CC BY 4.0)}
\begin{document}
\maketitle

\noindent\emph{Source and licence.} LLM-Text Watermarking based on Lagrange Interpolation. Version arXiv:2505.05712v4 (CC BY 4.0), distributed by arXiv under the Creative Commons Attribution 4.0 International licence, http://creativecommons.org/licenses/by/4.0/, as recorded in the arXiv licence metadata for this exact version and shown on the abstract page https://arxiv.org/abs/2505.05712v4. Full licence notice: \texttt{licenses/arxiv-2505.05712-CC-BY-4.0.txt}.\par\smallskip

\noindent\textit{Editorial scope.} Counted unit: the article's theorem thm:rand\_collinearity (source0.tex lines 367--377). The article's complete original proof of it is delivered with it (source0.tex lines 551--611, one \texttt{proof} environment of 61 lines, which the article places away from the statement). No mathematics is written, repaired or re-derived: the statement and the proof are the article's own lines, in the article's own notation, and no retained line is substituted or re-flowed.  No further source-local context is needed: every cross-reference of every retained line resolves inside the two delivered spans.  Nothing was omitted from any retained span, so the package carries no omission record.  No retained line contains a citation, so the wrapper carries no bibliography and no bibliography entry.  No retained line contains a figure environment, an image-inclusion command or drawing code, so no image is delivered and the package declares no asset dependency.  Editorial formatting only, in addition: an AMS \texttt{amsart} preamble that restates the article's own declarations and macros used by the retained lines, namely the environments \texttt{theorem} on separate counters, together with the standard packages those lines need.  The retained environments print in this document as Theorem 1 in order of appearance, whereas the article prints the same items with the numbering of its own full document.  The complete native source of the article as distributed in the arXiv e-print for this exact version is bundled unchanged as \texttt{sources/177836/source0.tex}.
\begin{theorem}[Random collinearity in $\mathrm{GF}(2^n)^2$]\label{thm:rand_collinearity}
Let $q = 2^n$. Let $S=\{p_1,\ldots,p_R\}$ be $R$ points sampled i.i.d.\ uniformly from $\mathrm{GF}(q)\times\mathrm{GF}(q)$, and consider only non-vertical lines of the form $y=ax+b$ with $a,b\in\mathrm{GF}(q)$. For any integer $k\ge 3$, the probability that there exists a non-vertical line containing at least $k$ points of $S$ satisfies

\[
    \Pr\big[\exists\, k\text{-collinear points}\big]
    \;\le\;
    \binom{R}{k}\, q^{-(k-2)}.
\]

A proof of the theorem is given in Appendix.
\end{theorem}

\begin{proof}[Proof of Theorem~\ref{thm:rand_collinearity}]
Let $q=2^n$ and let $S=\{p_1,\ldots,p_R\}$ be $R$ points sampled i.i.d.\ uniformly from $\GF(q)\times\GF(q)$. We consider only non-vertical lines of the form $y=ax+b$ with $a,b\in\GF(q)$.

\paragraph{Counting non-vertical lines.}
Each ordered pair $(a,b)\in\GF(q)\times\GF(q)$ determines a unique non-vertical line $\ell_{a,b}=\{(x,y): y=ax+b\}$, and distinct pairs give distinct lines. Hence the set $\mathcal{L}$ of non-vertical lines has cardinality $|\mathcal{L}|=q^2$.

\paragraph{Membership probability and the per-line distribution.}
Fix any $\ell\in\cL$. For a single random point $X$ uniformly distributed on $\GF(q)^2$, we have

\[
    \Pr[X\in \ell] \;=\; \frac{|\ell|}{|\GF(q)^2|} \;=\; \frac{q}{q^2} \;=\; \frac{1}{q}.
\]

Because $p_1,\ldots,p_R$ are sampled i.i.d., the random variable

\[
    N_\ell \;\defeq\; |\ell \cap S|
\]

is binomial: $N_\ell \sim \mathrm{Binomial}(R,1/q)$. In particular,

\[
    \Pr\!\big[N_\ell = k\big] \;=\; \binom{R}{k}\left(\frac{1}{q}\right)^k\!\left(1-\frac{1}{q}\right)^{R-k},
\]

which proves the ``per-line'' part of the theorem.

\paragraph{Existence of a $k$-collinear subset on some non-vertical line (union bound)}
For a line $\ell\in\cL$ and a $k$-subset $T\subseteq [R]$ (indexing the points in $S$), let

\[
    E_{\ell,T}\;=\;\{\text{all points } \{p_i\}_{i\in T} \text{ lie on } \ell\}.
\]

Then $\Pr[E_{\ell,T}] = (1/q)^k$, since the $p_i$ are independent and each falls on $\ell$ with probability $1/q$. The event that there exist at least $k$ collinear points on some non-vertical line is contained in

\[
    \bigcup_{\ell\in\cL}\;\bigcup_{T\in \binom{[R]}{k}} E_{\ell,T}.
\]

Applying the union bound,

\[
    \Pr\!\big[\exists\, k\text{-collinear points on a non-vertical line}\big]
    \;\le\;
\]
\noindent
\[
    \sum_{\ell\in\cL}\;\sum_{T\in \binom{[R]}{k}} \Pr[E_{\ell,T}]
    \;=\;
    q^2\cdot \binom{R}{k}\cdot \left(\frac{1}{q}\right)^{k}
    \;=\;
\]
\noindent
\[
    \binom{R}{k}\, q^{-(k-2)}.
\]

This is exactly the claimed bound.

\end{proof}

\end{document}
